% Source coverage: physical PDF pp. 25--30 (printed pp. 2--7), ending before Section 15.2.
% Equations: (15.1.1)--(15.1.19), plus intervening unnumbered displays.
% Footnotes: two (physical pp. 25 and 27). Divider: physical p. 29. Uncertainties: none.

\subsection{Gauge Invariance}
\label{sec:15.1}

We assume that the Lagrangian of our theory is invariant under a set of
infinitesimal transformations on the matter fields $\psi_\ell(x)$
\begin{equation}
 \delta\psi_\ell(x)=\ii\epsilon^a(x)(t_a)_\ell{}^m\psi_m(x) ,
 \tag{15.1.1}\label{eq:15.1.1}
\end{equation}
with some set of independent constant matrices\footnote{We label internal
symmetry generators with lowercase Latin letters $a$, $b$, etc.  Later, when
broken and unbroken generators occur together, letters $a$, $b$, etc. label
generators of spontaneously broken symmetries, while $i$, $j$, etc. label
generators of unbroken symmetries. A later lowercase Latin run is used for a
second internal family whenever needed to avoid a collision.}
$t_a$, and with real infinitesimal parameters $\epsilon^a(x)$ which (as
for gauge transformations in electrodynamics) are allowed to depend on position
in spacetime. We assume that these symmetry transformations are the infinitesimal
part of a Lie group; as shown in Section 2.2, this requires that the $t_a$
obey commutation relations
\begin{equation}
 [t_a,t_b]=\ii C^c{}_{ab}t_c ,
 \tag{15.1.2}\label{eq:15.1.2}
\end{equation}
where $C^c{}_{ab}$ are a set of real constants, known as the
structure constants of the group. The antisymmetry of the commutator immediately
tells us that the structure constants are similarly antisymmetric:
\begin{equation}
 C^c{}_{ab}=-C^c{}_{ba} .
 \tag{15.1.3}\label{eq:15.1.3}
\end{equation}
Also, from the Jacobi identity
\begin{equation}
 0=[[t_a,t_b],t_c]+[[t_c,t_a],t_b]
   +[[t_b,t_c],t_a]
 \tag{15.1.4}\label{eq:15.1.4}
\end{equation}
we see that the $C$s satisfy the further constraint
\begin{equation}
 0=C^d{}_{ab}C^e{}_{\dd c}
  +C^d{}_{ca}C^e{}_{\dd b}
  +C^d{}_{bc}C^e{}_{\dd a} .
 \tag{15.1.5}\label{eq:15.1.5}
\end{equation}
Any set of constants $C^c{}_{ab}$ that satisfy
Eqs.~\eqref{eq:15.1.3} and \eqref{eq:15.1.5} define at least one set of matrices
$t^A_a$:
\begin{equation}
 (t^A_a)^b{}_c=-\ii C^b{}_{ca} ,
 \tag{15.1.6}\label{eq:15.1.6}
\end{equation}
that satisfy the commutation relations \eqref{eq:15.1.2} with structure constants
$C^c{}_{ab}$:
\begin{equation}
 [t^A_a,t^A_b]=\ii C^c{}_{ab}t^A_c .
 \tag{15.1.7}\label{eq:15.1.7}
\end{equation}
This is known as the `adjoint' (or `regular') representation of the Lie algebra
with structure constants $C^c{}_{ab}$.

For example, in the original Yang--Mills theory, the matter fields were the
doublet consisting of proton and neutron fields $\psi_p$ and $\psi_n$:
\[
 \psi=\begin{pmatrix}\psi_p\\ \psi_n\end{pmatrix}
\]
and the $t_a$ with $a=1,2,3$ were the isospin matrices
\[
 t_1=\frac12\begin{pmatrix}0&1\\1&0\end{pmatrix},\qquad
 t_2=\frac12\begin{pmatrix}0&-\ii\\\ii&0\end{pmatrix},\qquad
 t_3=\frac12\begin{pmatrix}1&0\\0&-1\end{pmatrix} .
\]
These satisfy the commutation relations \eqref{eq:15.1.2} with
\[
 C^c{}_{ab}=\epsilon_{cab} ,
\]
where as usual $\epsilon_{cab}$ is $+1$ or $-1$ if
$c,a,b$ is an even or an odd permutation of 1, 2, 3,
respectively, and vanishes otherwise. We recognize this as the same as the Lie
algebra (2.4.18) of the three-dimensional rotation group; the matrices
$t_a$ here furnish what we recognize as the spin $1/2$ representation of
this Lie algebra. The matrices \eqref{eq:15.1.6} of the adjoint representation
are here (in a basis with rows and columns labelled 1,2,3):
\[
 t^A_1=\begin{pmatrix}0&0&0\\0&0&-\ii\\0&\ii&0\end{pmatrix},\qquad
 t^A_2=\begin{pmatrix}0&0&\ii\\0&0&0\\-\ii&0&0\end{pmatrix},\qquad
 t^A_3=\begin{pmatrix}0&-\ii&0\\\ii&0&0\\0&0&0\end{pmatrix} .
\]
This is the spin 1 representation of the Lie algebra of the rotation group.

Now consider what is needed to make the Lagrangian invariant under the
transformations \eqref{eq:15.1.1}. If there were no derivatives acting on the
fields, the task would be easy---any function of the matter fields that was
invariant under the transformation \eqref{eq:15.1.1} with $\epsilon^a$
constant would also be invariant with $\epsilon^a$ arbitrary real functions
of the spacetime coordinates. This is not the case if the Lagrangian involves
derivatives of the fields (as it must), because with position-dependent functions
$\epsilon^a(x)$, the derivatives of the matter fields do not transform like
the fields themselves. Differentiating Eq.~\eqref{eq:15.1.1} gives
\begin{equation}
 \delta\bigl(\partial_\mu\psi_\ell(x)\bigr)
 =\ii\epsilon^a(x)(t_a)_\ell{}^m\bigl(\partial_\mu\psi_m(x)\bigr)
 +\ii\bigl(\partial_\mu\epsilon^a(x)\bigr)(t_a)_\ell{}^m\psi_m(x) .
 \tag{15.1.8}\label{eq:15.1.8}
\end{equation}
To make the Lagrangian invariant, we need a field $A^a_\mu$, whose
transformation rule involves a term $\partial_\mu\epsilon^a$, which can be
used to cancel the second term in Eq.~\eqref{eq:15.1.8}. Since this field carries
an $a$-index, we would expect it also to undergo a matrix transformation
like Eq.~\eqref{eq:15.1.1}, but with $t_a$ replaced with the adjoint
representation matrices \eqref{eq:15.1.6}. Let us therefore tentatively take the
transformation relation of these new `gauge' fields as
\[
 \delta A^b_\mu=\partial_\mu\epsilon^b
 +\ii\epsilon^a(t^A_a)^b{}_c A^c_\mu
\]
or, using Eq.~\eqref{eq:15.1.6},
\begin{equation}
 \delta A^b_\mu=\partial_\mu\epsilon^b
 +C^b{}_{ca}\epsilon^a A^c_\mu .
 \tag{15.1.9}\label{eq:15.1.9}
\end{equation}
This allows us to construct a `covariant derivative':\footnote{As discussed in
the next section, in writing Eq.~\eqref{eq:15.1.10} we are tacitly supposing that
any coupling-constant factors like the electric charge are included in the
$t_b$, and hence also in the structure constants.}
\begin{equation}
 (D_\mu\psi(x))_\ell=\partial_\mu\psi_\ell(x)
 -\ii A^b_\mu(x)(t_b)_\ell{}^m\psi_m(x) .
 \tag{15.1.10}\label{eq:15.1.10}
\end{equation}
As planned, the term $\partial_\mu\epsilon^b$ in the transformation of
$A^b_\mu$ in the second term of Eq.~\eqref{eq:15.1.10} cancels the term
proportional to $\partial_\mu\epsilon^b$ in the transformation of the first
term, leaving us with
\begin{align*}
 \delta(D_\mu\psi)_\ell={}&\ii\epsilon^a(t_a)_\ell{}^m
 \partial_\mu\psi_m-\ii C^b{}_{ca}\epsilon^a
 A^c_\mu(t_b)_\ell{}^m\psi_m\\
 &{}+A^c_\mu(t_c)_\ell{}^m
 \ii\epsilon^a(t_a)_m{}^n\psi_n
\end{align*}
or, using Eq.~\eqref{eq:15.1.2},
\begin{equation}
 \delta(D_\mu\psi)_\ell=\ii\epsilon^a(t_a)_\ell{}^m(D_\mu\psi)_m ,
 \tag{15.1.11}\label{eq:15.1.11}
\end{equation}
so that $D_\mu\psi$ transforms just like $\psi$ itself.

We also need to worry about derivatives of the gauge field. In order to eliminate
the term $\partial_\nu\partial_\mu\epsilon^b$ in the transformation of
$\partial_\nu A^b_\mu$, we antisymmetrize with respect to $\mu$ and $\nu$,
just as in electrodynamics. However, we still have terms in the transformation of
$\partial_\nu A^b_\mu-\partial_\mu A^b_\nu$ proportional to first
derivatives of $\epsilon(x)$, arising from the second term in
Eq.~\eqref{eq:15.1.9}. The easiest way to construct a `covariant curl',
$F^c_{\nu\mu}$, in whose transformation rule all such derivatives of
$\epsilon(x)$ cancel, is to consider the commutator of two covariant derivatives
acting on a matter field $\psi$:
\begin{equation}
 \bigl([D_\nu,D_\mu]\psi\bigr)_\ell
 =-\ii(t_c)_\ell{}^m F^c_{\nu\mu}\psi_m ,
 \tag{15.1.12}\label{eq:15.1.12}
\end{equation}
where
\begin{equation}
 F^c_{\nu\mu}\equiv\partial_\nu A^c_\mu
 -\partial_\mu A^c_\nu
 +C^c{}_{ab}A^a_\nu A^b_\mu .
 \tag{15.1.13}\label{eq:15.1.13}
\end{equation}
Eq.~\eqref{eq:15.1.12} makes it obvious that $F^c_{\nu\mu}$ must
transform just like a matter field that happens to belong to the adjoint
representation:
\begin{equation}
 \delta F^b_{\nu\mu}
 =\ii\epsilon^a(t^A_a)^b{}_c F^c_{\nu\mu}
 =\epsilon^a C^b{}_{ca}F^c_{\nu\mu} .
 \tag{15.1.14}\label{eq:15.1.14}
\end{equation}
The reader may check by direct calculation (using the relation
\eqref{eq:15.1.5}) that the quantity $F^a_{\nu\mu}$ defined in
\eqref{eq:15.1.13} actually has the simple transformation rule
\eqref{eq:15.1.14}.

For some purposes, it is useful to know that these infinitesimal gauge
transformations can be upgraded to finite transformations. A group element can
be parameterized by a set of real functions $\Lambda^a(x)$ so that it acts
on a general matter field $\psi_\ell(x)$ through the matrix transformation
\begin{equation}
 \psi_\ell(x)\longrightarrow\psi_{\ell\Lambda}(x)
 =\bigl[\exp(\ii t_a\Lambda^a(x))\bigr]_\ell{}^m\psi_m(x) .
 \tag{15.1.15}\label{eq:15.1.15}
\end{equation}
We want the covariant derivative to transform in the same way:
\begin{equation}
 (\partial_\mu-\ii t_a A^a_{\mu\Lambda})\psi_\Lambda
 =\exp(\ii t_a\Lambda^a)(\partial_\mu-\ii t_a A^a_\mu)\psi ,
 \tag{15.1.16}\label{eq:15.1.16}
\end{equation}
so we must impose on $A^a_\mu$ the transformation rule
$A^a_\mu\longrightarrow A^a_{\mu\Lambda}$, with
\[
 \partial_\mu\exp(\ii t_b\Lambda^b)
 -\ii t_b\exp(\ii t_a\Lambda^a)A^b_{\mu\Lambda}
 =-\ii\exp(\ii t_a\Lambda^a)t_b A^b_\mu
\]
or in other words
\begin{align}
 t_a A^a_{\mu\Lambda}={}&
 \exp(\ii t_b\Lambda^b)t_a A^a_\mu
 \exp(-\ii t_b\Lambda^b)\notag\\
 &{}-\ii\bigl[\partial_\mu\exp(\ii t_b\Lambda^b)\bigr]
 \exp(-\ii t_b\Lambda^b) .
 \tag{15.1.17}\label{eq:15.1.17}
\end{align}
Eqs.~\eqref{eq:15.1.15} and \eqref{eq:15.1.17} reduce to the previous
transformation rules \eqref{eq:15.1.1} and \eqref{eq:15.1.9} in the limit where
$\Lambda^a(x)$ is an infinitesimal $\epsilon^a(x)$.

From Eq.~\eqref{eq:15.1.17}, we can see that by a suitable choice of
$\Lambda^b(x)$, it is always possible to make $A^a_{\mu\Lambda}(x)$
vanish at any one point, say $x=z$. (Simply take $\Lambda^a(z)$ to vanish,
and $\partial\Lambda^a(x)/\partial x^\mu=-A^a_\mu(x)$ at $x=z$.)
Also, it is always possible to choose $\Lambda^b(x)$ so that any one
spacetime component of $A^a_{\mu\Lambda}(x)$ vanishes for all $a$
everywhere in at least a finite domain around any given point. For instance, to
make $A^a_{3\Lambda}(x)$ vanish, we must solve the set of ordinary
first-order differential equations for the parameters $\Lambda^b(x)$:
\begin{equation}
 \partial_3\exp(\ii t_b\Lambda^b)
 =-\ii\exp(\ii t_b\Lambda^b)t_a A^a_3 ,
 \tag{15.1.18}\label{eq:15.1.18}
\end{equation}
which always have a solution in at least a finite domain around any ordinary
point.

However, in general it is not possible to choose $\Lambda^a(x)$ to make all
four components $A^a_{\mu\Lambda}(x)$ vanish in a finite region. For this
purpose, we would have to be able to satisfy the partial differential equations
\begin{equation}
 \partial_\mu\exp(\ii t_b\Lambda^b)
 =-\ii\exp(\ii t_b\Lambda^b)t_a A^a_\mu ,
 \tag{15.1.19}\label{eq:15.1.19}
\end{equation}
which cannot be solved unless certain integrability conditions are satisfied. In
particular, if $A^a_{\mu\Lambda}(x)$ vanishes everywhere then so does
$F^a_{\mu\nu\Lambda}(x)$, but since the field strength transforms
homogeneously, $F^a_{\mu\nu\Lambda}(x)$ can vanish only if
$F^a_{\mu\nu}(x)$ does. A gauge field $A^a_\mu(x)$ is called a
`pure gauge' field if there exists a gauge transformation which makes it vanish
everywhere. It is not difficult to show that the condition that
$F^a_{\mu\nu}$ should vanish everywhere is not only necessary but also
sufficient for $A^a_\mu(x)$ to be expressible in any simply connected
region as a pure gauge field~\hyperref[ch15-ref-6]{[6]}.

\begin{center}
$*\qquad *\qquad *$
\end{center}

There is a deep analogy between the construction here of objects that transform
simply under gauge transformations and the construction in general relativity of
objects that transform covariantly under general coordinate transformations. Just
as we use the gauge field to construct covariant derivatives $D_\mu\psi_\ell$ of
matter fields with the same gauge transformation properties as the matter fields
themselves, so we use the affine connection $\Gamma^\mu{}_{\nu\lambda}(x)$ to
construct covariant derivatives of tensors $T^{\rho\cdots}{}_{\kappa\lambda\cdots}$:
\[
 T^{\rho\cdots}{}_{\kappa\cdots;\nu}
 \equiv\partial_\nu T^{\rho\cdots}{}_{\kappa\cdots}
 +\Gamma^\rho{}_{\nu\lambda}T^{\lambda\cdots}{}_{\kappa\cdots}+\cdots
 -\Gamma^\mu{}_{\nu\kappa}T^{\rho\cdots}{}_{\mu\cdots}-\cdots ,
\]
which are themselves tensors. Also, from the derivatives of the gauge field we
constructed a field strength $F^a_{\mu\nu}$ with the gauge transformation
property of a matter field belonging to the adjoint representation of the gauge
group; correspondingly, from the derivatives of the affine connection we may
construct a quantity:
\[
 R^\lambda{}_{\mu\nu\kappa}
 =\frac{\partial\Gamma^\lambda{}_{\mu\nu}}{\partial x^\kappa}
 -\frac{\partial\Gamma^\lambda{}_{\mu\kappa}}{\partial x^\nu}
 +\Gamma^\eta{}_{\mu\nu}\Gamma^\lambda{}_{\kappa\eta}
 -\Gamma^\eta{}_{\mu\kappa}\Gamma^\lambda{}_{\nu\eta} ,
\]
which transforms as a tensor, the Riemann--Christoffel curvature tensor. The
commutator of two gauge-covariant derivatives $D_\mu$ and $D_\nu$ may be
expressed in terms of the field-strength tensor $F^a_{\mu\nu}$; similarly,
the commutator of two covariant derivatives with respect to $x^\nu$ and $x^\kappa$
may be expressed in terms of the curvature:
\[
 T^{\lambda\cdots}{}_{\mu\cdots;\nu;\kappa}
 -T^{\lambda\cdots}{}_{\mu\cdots;\kappa;\nu}
 =R^\lambda{}_{\sigma\nu\kappa}T^{\sigma\cdots}{}_{\mu\cdots}+\cdots
 -R^\sigma{}_{\mu\nu\kappa}T^{\lambda\cdots}{}_{\sigma\cdots}-\cdots .
\]
The necessary and sufficient condition for the existence of a gauge in which the
gauge field vanishes in a finite simply connected region is the vanishing of the
field-strength tensor, and the necessary and sufficient condition for the
existence of a coordinate system in which the affine connection vanishes in a
finite simply connected region is the vanishing of the Riemann--Christoffel
curvature tensor. The analogy breaks down in one important respect: in general
relativity the affine connection is itself constructed from first derivatives of
the metric tensor, while in gauge theories the gauge fields are not expressed in
terms of any more fundamental fields.
